% 2026-11-11 / v31
\section{Artin基底上のJacobian sectorと導来自然変換の底変更}
\label{ras:sec:artin-sector}

本節では標数零の体$k$、$R_m=k[h]/(h^{m+1})$、
$S_m=R_m[x_1,\ldots,x_5]$、次数三の
$W_m\in(S_m)_3$を考える。$F=W_m\bmod h$は
$\PP^4_k$の滑らかなcubic threefoldを定めるとする。
目的は、相対的な五変数のJacobian sectorと、graded MF圏の
$\Id\to(1),(2)$の導来自然変換を区別しながら底変更を調べることである。
以下の$\RHom$は常に\emph{$R_m$-線形}な導来Homである。
$R_m$の絶対的なHochschild理論を取り込まない。

体上のkernelの計算にはBallard--Favero--Katzarkov\cite{RASBFK}、
Lin--Zhang\cite[Theorem~3.5, Remark~3.6]{RASLZ}を用いる。
残余体のstabilizationが生成元となる体上の結果は
Dyckerhoff\cite{RASDyck}を用いる。
以下では、非正則Artin底へ適用するときに必要な
平坦複体の持上げと、可換対角の有限模型を明示する。
これら既知の結果そのものの新規性は主張しない。

\subsection{勾配Koszul複体の相対完全性}

\begin{lemma}[勾配の全Artin完全性]\label{ras:lem:gradient}
$K_m=K_\bullet(\partial_1W_m,\ldots,\partial_5W_m;S_m)$を
通常のKoszul鎖複体とする。このとき
\[
 H_i(K_m)=0\quad(i>0),\qquad
 H_0(K_m)=J_m:=S_m/(\partial_1W_m,\ldots,\partial_5W_m).
\]
$J_m$は内部次数別に有限自由な$R_m$-加群で、Hilbert級数は
\[
 \sum_{j\ge0}\operatorname{rank}_{R_m}(J_m)_jt^j=(1+t)^5.
 \tag{RAS1}\label{ras:eq:hilbert}
\]
特に$(J_m)_1\simeq R_m^5$、$(J_m)_2\simeq R_m^{10}$である。
全ての商$R_m\to R_a$（$a\le m$）への底変更は
$J_m\otimes_{R_m}R_a\simeq J_a$と両立する。
\end{lemma}
\begin{proof}
$F$の偏微分は原点以外に共通零点を持たない。
$k[x_1,\ldots,x_5]$はCohen--Macaulayなので、五つの二次式は
正則列であり、$K_0$は次数零以外で完全である。
従って$J_0$のHilbert級数は
$(1-t^2)^5/(1-t)^5=(1+t)^5$である。

$m\ge1$のとき、各成分は$R_m$-自由なので
\[
 0\longrightarrow K_0\xrightarrow{h^m}K_m
 \longrightarrow K_{m-1}\longrightarrow0
 \tag{RAS2}\label{ras:eq:short}
\]
は鎖複体の短完全列である。左端の微分は
$h^m\partial W_m=h^m\partial F$なので、左端は実際に$K_0$である。
長完全列による帰納法で$H_i(K_m)=0$（$i>0$）を得て、
\[
 0\longrightarrow J_0\xrightarrow{h^m}J_m
 \longrightarrow J_{m-1}\longrightarrow0
\]
も得る。各内部次数$j$の$J_m$の$k$-長さは
$(m+1)\dim_k(J_0)_j$である。
$(J_0)_j$の基底を同次数の多項式に持ち上げると、Nakayamaにより
$R_m^{d_j}\twoheadrightarrow(J_m)_j$を得る。
両辺の$k$-長さは$(m+1)d_j$なので、これは同型である。
同じ持上げ基底を使えば任意の$R_a$への底変更が従う。
\end{proof}

\begin{lemma}[有限nilpotent比較の検出]\label{ras:lem:nilpotent}
$f:C^\bullet\to D^\bullet$を、各項が$R_m$-平坦な複体の射とする。
$f\otimes_{R_m}k$が擬同型なら、$f$は擬同型である。
有界性は仮定しないが、\emph{各項の平坦性}は必要である。
\end{lemma}
\begin{proof}
$E=\operatorname{Cone}(f)$も項別平坦である。
有限フィルトレーション
$E\supset hE\supset\cdots\supset h^mE\supset0$の各商は、
平坦性により$E/hE$と複体として同型である。
$E/hE$は非輪状であるから、有限回の短完全列を用いて$E$も
非輪状となる。
\end{proof}

\begin{example}[特殊ファイバーの次元だけでは足りない]
\label{ras:ex:basechange-fails}
$R_1=k[h]/(h^2)$上の二項複体
$C=[R_1\xrightarrow{h}R_1]$を次数$0,1$に置く。
$H^0(C)=hR_1\simeq k$だが、
$H^0(C)\otimes_{R_1}k\to H^0(C\otimes_{R_1}k)=k$は\emph{零写像}である。
両コホモロジーの$k$-次元が一であっても底変更の同型は従わない。
従って、後述の標的の底変更を体上の次元公式だけで済ませてはならない。
\end{example}

\subsection{体上のkernel比較をArtin底へ持ち上げる際の構成}

$(\mathcal T_m,(1))$を、有限rank自由な次数付き
matrix factorizationとその冪等完備化から作った$R_m$-線形dg圏とする。
そのdg双加群による自然変換を
\[
 T_{r,m}:=H^0\RHom_{\mathcal T_m^e}
 (\Delta_{\mathcal T_m},\Delta_{\mathcal T_m}(r))
 \qquad(r=1,2)
 \tag{RAS3}\label{ras:eq:target}
\]
とする。

\begin{lemma}[導来bar評価の条件付き持上げ]\label{ras:lem:generator}
$\mathcal C_m$を項別$R_m$-平坦なHomを持つdg圏、
$G_m\in\mathcal C_m$とする。各対象$X,Y$について、
$G_m$を介した\emph{導来二側bar評価}の鎖模型があり、
その始域と終域が項別$R_m$-平坦で、$h=0$への剰余が
体上のMorita評価と厳密に一致すると仮定する。
体上で$G_0$がMorita生成元なら、この評価は擬同型である。
従って$G_m$は、この導来bar模型でのMorita生成元となる。
\end{lemma}
\begin{proof}
特殊ファイバーの評価はMorita生成性から擬同型である。
与えられた導来bar評価のconeは項別平坦で、
その剰余は特殊ファイバーで非輪状だから
補題\ref{ras:lem:nilpotent}によりcone自身も非輪状である。
\end{proof}
\begin{remark}[一般の特殊生成子持上げは偽]
指定した導来bar評価の存在・導来性・特殊化整合性は
単なるHom複体の項別平坦性からは従わない。
具体的な反例と隣接sectorの障害は
後続のArtin障害の節に記した。
\end{remark}

五変数の場合、$W_m=\sum_i x_i(\partial_iW_m/3)$から
残余点のstabilization $K_m$をKoszul型の有限rank因子分解として
作れる。$K_m,(1)K_m,(2)K_m$の直和を$G_m$とする。
体上ではこれがgraded MFのMorita生成元となる\cite{RASDyck,RASBFK}。
しかしこのことだけからArtin底の生成性は従わない。
補題\ref{ras:lem:generator}の導来bar評価・剰余整合性を別途確認する。
上のbarの各項は同次多項式の有限自由$R_m$-加群の直和である。
自然変換を計算するbar \emph{cochain}では直積が現れるが、
$R_m$はNoetherian（従ってcoherent）であり、自由加群の直積は平坦である。
$R_m\to k$への剰余は、各係数を成分毎に剰余する操作と可換する。

差$W_m(x)-W_m(y)$を
\[
 W_m(x)-W_m(y)=\sum_i(x_i-y_i)D_iW_m(x,y)
\]
と分解し、$\sum_i(x_i-y_i)\iota_i+\sum_iD_iW_m\,\wedge e_i$を
微分とするKoszul因子分解$\mathscr D_m$を取る。
これは$\mathbb G_m$-同変な\emph{stabilized diagonal kernel}である。
$x_i-y_i$は任意の$R_m$上で正則列なので、有限rank因子分解への畳込みにより
恒等関手を与えることは、通常のKoszul contractionで確認できる。
内部shiftを持つ$\mathscr D_m(r)$も同様である。

\begin{proposition}[相対kernel作用のnilpotent検出]\label{ras:prop:kernel-action}
$\mathscr D_m$によるkernel作用から得られる、相対的な自然な複体の射
\[
 \kappa_{r,m}:\Hom_{\mathrm{MF}^{\rm gr}(W_m(x)-W_m(y))}
 (\mathscr D_m,\mathscr D_m(r))
 \longrightarrow
 \RHom_{\mathcal T_m^e}(\Delta,\Delta(r))
 \tag{RAS5}\label{ras:eq:kappa}
\]
を、右辺で$G_m$の二側bar模型を使って表す。
この比較の特殊ファイバーが体上のkernel定理による擬同型なら、
$\kappa_{r,m}$は全ての$m$で擬同型である。
\end{proposition}
\begin{proof}
左辺は有限自由因子分解の同次Homなので項別$R_m$-自由である。
右辺の$G_m$-bar模型の各項は上記のように$R_m$-平坦で、
特殊化は体上の同じbar模型となる。kernelの作用は
$R_m$上の畳込み・合成で作るため、特殊化と可換である。
従って特殊化が擬同型なら、補題\ref{ras:lem:nilpotent}を適用できる。
\end{proof}

\begin{remark}[比較の適用範囲]
体上のkernel定理\cite{RASBFK}は、そのdg enhancementと
連続kernel functorを指定する。式\eqref{ras:eq:kappa}について
\emph{同じ}enhancementを用い、相対畳込みと$G_m$-bar評価を
整合させる必要がある。異なる第二種Hochschild totalizationを
同じ名前だからと同一視してはならない。
以下の計算は、有限自由のstabilized diagonalの相対模型について行う。
\end{remark}

\subsection{恒等sectorと非恒等sector}

\begin{proposition}[五変数の相対sector計算]\label{ras:prop:sectors}
標準の$\mu_3$-sector分解を相対stabilized diagonalに適用すると、
$\mathscr D_m$の内部shift付き自己準同型の次数零は
\[
 H^0\Hom(\mathscr D_m,\mathscr D_m(r))
 \simeq (J_m)_r\quad(r=1,2)
 \tag{RAS6}\label{ras:eq:sectors}
\]
となり、積は$J_m$の通常の多項式積である。
この同型は全$R_m\to R_a$と両立する。
\end{proposition}
\begin{proof}
$\mu_3$は標数零で線型簡約だから、sectorへの射影は
$R_m$上の直和分解である。恒等sectorでは、二つの対角の
Koszul homotopyを$y=x$で縮約すると、偏微分
$(\partial_1W_m,\ldots,\partial_5W_m)$のKoszul複体を得る。
補題\ref{ras:lem:gradient}により、次数零は$J_m$で、
正のKoszulホモロジーは零である。

非恒等$\zeta\in\mu_3$では、twisted対角との交点に
$((1-\zeta)x_1,\ldots,(1-\zeta)x_5)$のKoszul複体が現れる。
$1-\zeta$は単元なので、この五つは$S_m$上の正則列である。
その縮約後に残るnormal determinantは\emph{奇数}（階数五）の
cohomological parityを持つ。$(3)=[2]$によるシフトは偶数なので、
どの内部shift$r$でもcohomological degree零には寄与しない。
この計算は体上のsector formula\cite[Theorem~3.5]{RASLZ}の
明示的Koszul証明を$R_m$上で行ったものであり、
$1-\zeta$の可逆性以外に体の性質を使わない。
恒等sectorの積はdiagonalへの多項式倍であり、積の同定も従う
\cite[Remark~3.6]{RASLZ}。
\end{proof}

\begin{theorem}[相対kernel比較の下での(BC)]\label{ras:thm:bc}
各$R_m$について、\emph{指定生成子の導来二側bar評価が
補題\ref{ras:lem:generator}の条件を満たし、かつ相対kernel作用の射
\eqref{ras:eq:kappa}が体上のBFK kernel比較と同じenhancementにおいて
定義され、その$h=0$剰余が擬同型である}とする。
このとき、論文Bの条件(BC)がすべての$m$について成り立つ。
特に
\[
 T_{1,m}\simeq R_m^5,\quad T_{2,m}\simeq R_m^{10},\quad
 T_{r,m}\otimes_{R_m}k\xrightarrow{\sim}T_{r,0}\quad(r=1,2)
\]
であり、低次数の積は形式可換模型$W_m$のJacobian積である。
\end{theorem}
\begin{proof}
補題\ref{ras:lem:generator}の追加仮定により、導来bar模型で
実際の導来自然変換を計算できる。
命題\ref{ras:prop:kernel-action}のnilpotent持上げと
命題\ref{ras:prop:sectors}を合成すれば、標的の同定と底変更を得る。
内部shiftはdiagonalのequivariant characterのshiftであり、
比較の畳込みに組み込まれている。積の比較はkernelの合成であり、
恒等sectorでは多項式積となる。
\end{proof}

\begin{corollary}[混合族に関する条件付き正式]\label{ras:cor:formal-G}
既存の有限Artin MF同値・線束残余比較と、
上の相対kernel作用の仮定が与えられているとする。
$\lambda=1$で有限曲率付き対角から計算した多項式候補
$G_h+v^3$と、既存のformalityで選んだ可換形式三次式$W_h$は、
$\CC[[h]]$上、恒等に近い形式線形座標変換で同型である。
非零複素数$h$での圏同値はこの結論に含まれない。
\end{corollary}
\begin{proof}
定理\ref{ras:thm:bc}と既存の乗法的characteristic写像の
Artin判定から、全Artin次数で低次数積が一致する。
対応を逆極限に取り、次数二の核として勾配空間を同定する。
$h=0$では座標の同定を恒等に規格化できる。
有限代数的積分の定理とsymmetrizerの形式三乗根により、
両三次式を形式的な線形座標変換で同定できる。
\end{proof}

\subsection{残る検証の境界}

相対kernel作用の\emph{特殊ファイバー擬同型}は体上の既知定理だが、
非正則Artin上の有限自由enhancement、畳込み、二側barを
同じモデルで整合させた完全な鎖レベルの実装は、
この原稿では独立の有限行列証明書として提供していない。
従って上の定理は必要な比較モデルを指定した\emph{条件付き}の定理である。
前節の有限曲率付き対角と相対stabilized diagonalとの比較を
省略して、実際の非可換rotation積と無条件に同定してはならない。
最優先課題は、そのchain-level kernel作用の自然性・$R_m$-平坦性・
特殊化を一つの有限自由模型で具体的に検証することである。

具体的な五変数の古典特殊ファイバー
\[
 F_1=(x^3+y^3+z^3)/3+u^3+v^3+uxy
\]
では$F_x=x^2+uy$, $F_y=y^2+ux$, $F_u=3u^2+xy$から
$u\ne0$を仮定すると$(xy)^2=9u^4$かつ
$x^2y^2=-3u^4$となり矛盾する。従って唯一の勾配共通零点は原点。
厳密なGr\"obner証明書は
\texttt{computations/verify\_artin\_sectors.py}で再生成でき、
Hilbert関数は$(1,5,10,10,5,1)$である。
さらに計算済み候補$G_{1/2}+v^3$も同じ関数を持つことを検査する。
これは(BC)の鎖レベルの実装を検証したものではない。
